\chapter{Kiểm thử}\label{chap:testing}
\section{Mục tiêu và phạm vi}
Nhóm xây dựng kế hoạch kiểm thử để đánh giá các chức năng của iCareX và phát hiện những tình huống có thể làm sai dữ liệu. Trọng tâm gồm quyền truy cập theo tài khoản, tính đúng của chuỗi số km, cách tổng hợp chi phí và việc chuyển kỳ bảo dưỡng sau khi thay đổi lịch sử.

Kiểm thử đơn vị tập trung vào các quy tắc có thể đánh giá độc lập, như kiểm tra số km và tính hạn theo tháng. Kiểm thử API xem xét xác thực, CSRF, dữ liệu đầu vào và mã phản hồi. Kiểm thử tích hợp sử dụng PostgreSQL để kiểm tra quan hệ dữ liệu, giao dịch và cập nhật đồng thời. Kiểm thử toàn trình theo dõi thao tác từ trình duyệt đến dữ liệu đã lưu.

\section{Dữ liệu và môi trường kiểm thử}
\subsection{Bộ dữ liệu}
Nhóm sử dụng hai tài khoản A và B để kiểm tra phân quyền. Tài khoản A có chiếc xe trong kịch bản minh họa ở chương 1, với khoản nhiên liệu 200.000 VNĐ và dịch vụ 250.000 VNĐ. Tổng tiền cần đối chiếu là 450.000 VNĐ. Các trường hợp tính ngày dùng thời điểm kiểm thử cố định để kết quả có thể lặp lại.

Các trường hợp bổ sung gồm ngày cuối tháng, năm nhuận, ngày trong tương lai, số km bằng nhau, lịch sử nhập bổ sung và dữ liệu không hợp lệ. Mỗi kiểm thử tích hợp sử dụng dữ liệu riêng để tránh ảnh hưởng đến kết quả của trường hợp khác.

\subsection{Điều kiện thực hiện}
Cơ sở dữ liệu kiểm thử được tách biệt với dữ liệu đang sử dụng. Các bài kiểm tra cookie Secure được thực hiện qua HTTPS. Đối với kiểm thử hiệu năng, nhóm ghi nhận cấu hình máy, phiên bản phần mềm, số lượng bản ghi và tải đồng thời cùng kết quả đo.

\section{Các trường hợp kiểm thử}
Mỗi ca kiểm thử có mã TC, tiêu chí nghiệm thu liên quan, dữ liệu đầu vào, các bước thực hiện và kết quả mong đợi. Những ca dưới đây thuộc kế hoạch kiểm thử; kết quả thực tế được ghi nhận sau khi thực hiện, kèm trạng thái đạt, không đạt hoặc chưa thực hiện.
\begin{smalltable}\begin{longtable}{L{1.2cm}L{2.5cm}L{3.6cm}L{3.6cm}L{3.7cm}}
\caption{Kế hoạch kiểm thử: dữ liệu, bước thực hiện và kết quả mong đợi}\label{tab:testcases}\\
\toprule ID / AT & Mục tiêu & Tiền điều kiện / input & Bước chính & Kết quả mong đợi\\\midrule\endfirsthead
\toprule ID / AT & Mục tiêu & Tiền điều kiện / input & Bước chính & Kết quả mong đợi\\\midrule\endhead
TC01 /\allowbreak  AT01 & Đăng ký hợp lệ & Email mới; password đạt policy; CSRF đúng. & Gửi register; đọc DB; gọi /\allowbreak me trước login. & 201; User/\allowbreak Account/\allowbreak Settings cùng ID; chưa login nên /\allowbreak me 401.\\
TC02 /\allowbreak  AT01 & Email trùng sau chuẩn hóa & Đã có email account A. & Đăng ký lại cùng email khác casing/\allowbreak space. & 409; không thêm profile/\allowbreak settings thừa.\\
TC03 /\allowbreak  AT01 & Password yếu & Email mới,  password thiếu yêu cầu. & Gửi trực tiếp API bỏ qua UI validation. & 400 validation; không tạo account một phần.\\
TC04 /\allowbreak  AT01 & Login sai và lockout & Account tồn tại; dùng password sai. & Sai 5 lần; thử lại trong 15 phút; dùng controlled time nếu có. & Không tạo session; thông báo chung; lockout đúng policy.\\
TC05 /\allowbreak  AT01 & Logout revoke & A login có sid hợp lệ. & Logout +\allowbreak  CSRF; dùng lại cookie cũ gọi /\allowbreak me. & 204 rồi 401; DB revoked; cache A được xóa ở UI đích.\\
TC06 /\allowbreak  AT02 & GET/\allowbreak PUT/\allowbreak DELETE chéo owner & A có xe/\allowbreak reading; B login. & B dùng ID A gọi đọc/\allowbreak sửa/\allowbreak xóa. & 404/\allowbreak từ chối; không lộ field A; DB không đổi.\\
TC07 /\allowbreak  AT02 & Link plan chéo xe/\allowbreak owner & B có service; A có plan. & Gửi item B với planId A; thêm trường accountId A vào payload. & Không liên kết; server lấy owner từ session; rollback.\\
TC08 /\allowbreak  AT03 & Archive giữ dữ liệu & Xe active có hoạt động và plan. & Archive với version; đọc history; thử mọi đường ghi. & History còn; mutation bị chặn; nhắc chủ động dừng ở đích.\\
TC09 /\allowbreak  AT03 & Restore không reset kỳ & Xe archived đã qua hạn. & Restore với version; đọc plan. & Xe active; due theo mốc cũ và today hiện tại; không completion giả.\\
TC10 /\allowbreak  AT04 & ODO sai chuỗi quá khứ & 20/\allowbreak 10=10.200; 25/\allowbreak 10=10.300. & Chèn 22/\allowbreak 10=10.350,  rồi 10.250. & Lần đầu bị từ chối; lần sau lưu; current vẫn 10.300.\\
TC11 /\allowbreak  AT04 & ODO âm/\allowbreak ngày tương lai & todayVN=01/\allowbreak 11/\allowbreak 2026. & Gửi odo=-1; gửi date=02/\allowbreak 11/\allowbreak 2026. & Validation không lưu; ngày theo Việt Nam,  không theo host UTC.\\
TC12 /\allowbreak  AT04 & Cùng ngày và giá trị bằng nhau & Có reading 10.200 ngày 20/\allowbreak 10. & Thêm reading cùng ngày/\allowbreak ODO; query nhiều lần. & Cho phép theo đặc tả; thứ tự ổn định theo ODO và khóa phụ.\\
TC13 /\allowbreak  AT04 & Xóa max ODO & Manual readings 10.200 và 10.300. & Xóa max với version; query current. & Current về 10.200 hoặc mốc hợp lệ cao nhất còn lại.\\
TC14 /\allowbreak  AT04 & Đổi baseline phá chuỗi & Base 10.000,  history có 10.200. & Sửa base thành 10.250. & Từ chối,  baseline/\allowbreak history giữ nguyên.\\
TC15 /\allowbreak  AT05 & Fuel retry cùng payload & RequestId R1,  fuel=200.000. & Gửi 2 lần cùng key/\allowbreak payload,  kể cả mô phỏng mất response. & Chỉ 1 fuel/\allowbreak linked reading và 200.000 tiền; replay đúng contract.\\
TC16 /\allowbreak  AT05 & Fuel retry khác payload & R1 đã tạo 200.000. & Gửi R1 với amount=210.000. & 409; nguồn cũ không bị sửa,  không tạo nguồn mới.\\
TC17 /\allowbreak  AT05 & Sửa/\allowbreak xóa fuel nguồn & Fuel đã có linked reading. & Sửa ngày/\allowbreak ODO/\allowbreak tiền rồi xóa; refetch xe và report. & Nguồn/\allowbreak reading cùng thay đổi; không mồ côi; tổng phản ánh đúng.\\
TC18 /\allowbreak  AT06 & Service hai dòng & Repair=100.000,  Maintenance=150.000. & Tạo một service tại 10.300 với 2 items. & Một header,  hai items,  một reading; tổng 250.000.\\
TC19 /\allowbreak  AT06 & Rollback giữa transaction & Fixture service nhiều item; fault injection trước commit. & Gây lỗi sau insert header/\allowbreak reading hoặc khi đồng bộ nhắc. & Không còn header/\allowbreak item/\allowbreak reading/\allowbreak completion/\allowbreak key/\allowbreak nhắc một phần.\\
TC20 /\allowbreak  AT06 & Tổng tiền service không hợp lệ & Items âm hoặc mọi amount=0. & Gửi trực tiếp API với các biến thể. & Từ chối âm và tổng 0; item=0 cho phép nếu tổng dương.\\
TC21 /\allowbreak  AT07 & Biên Time và cuối tháng & Baseline 31/\allowbreak 01/\allowbreak 2026,  interval 1 tháng. & Evaluate với today=27/\allowbreak 02, 28/\allowbreak 02, 01/\allowbreak 03; ngưỡng rõ. & Due 28/\allowbreak 02; Soon/\allowbreak Due/\allowbreak Overdue theo biên. Năm 2028 due 29/\allowbreak 02.\\
TC22 /\allowbreak  AT07 & Distance bằng/\allowbreak vượt hạn & dueOdo=11.300; thresholdKm=100. & Evaluate current 11.199, 11.200, 11.300, 11.301. & NotDue,  Soon,  Due,  Overdue tương ứng.\\
TC23 /\allowbreak  AT07 & Either ưu tiên mức cao & Một chiều NotDue; chiều kia Due/\allowbreak Overdue. & Evaluate các tổ hợp đảo vị trí hai chiều. & Chọn mức cần chú ý cao nhất,  không đợi đủ cả hai.\\
TC24 /\allowbreak  AT07 & Either thiếu dữ liệu & Một chiều thiếu; chiều kia NotDue hoặc Soon. & Evaluate hai trường hợp và đọc reason. & NotDue+\allowbreak missing =$>$ Unknown; Soon+\allowbreak missing =$>$ Soon kèm flag.\\
TC25 /\allowbreak  AT07 & Baseline và ngưỡng & Mode Time/\allowbreak Distance/\allowbreak Either. & Gửi chu kỳ 0,  ngưỡng vượt kỳ,  baseline tương lai/\allowbreak vượt current. & Từ chối đúng field; không sinh service/\allowbreak chi phí giả.\\
TC26 /\allowbreak  AT08 & Completion chỉ plan liên kết & Hai plan cùng xe; service có Repair và Maintenance link plan 1. & Tạo service; đọc cả hai plan. & Chỉ plan 1 chuyển mốc; Repair không chuyển plan 2.\\
TC27 /\allowbreak  AT08 & Service nhập cũ không là latest & Đã có completion 25/\allowbreak 10. & Thêm completion 22/\allowbreak 10 hợp lệ sau baseline. & Latest vẫn 25/\allowbreak 10; kỳ không lùi theo thứ tự request.\\
TC28 /\allowbreak  AT08 & Xóa/\allowbreak sửa latest completion & Có completion 22/\allowbreak 10 và 25/\allowbreak 10. & Xóa 25/\allowbreak 10 rồi xóa 22/\allowbreak 10; sửa liên kết ở biến thể khác. & Mốc lần lượt 22/\allowbreak 10 rồi baseline; mọi plan cũ/\allowbreak mới được tính lại.\\
TC29 /\allowbreak  AT08 & Completion cùng ngày & Plan đã completion ngày 25/\allowbreak 10. & Thêm service khác cùng plan cùng ngày; thử request đồng thời. & Từ chối duplicate; transaction không để dữ liệu nguồn sai một phần.\\
TC30 /\allowbreak  AT09 & Worker lặp và restart & Một plan đang Soon/\allowbreak Due; reminders bật. & Chạy 2 vòng,  restart,  chạy lại; query khóa nhắc. & Tối đa một record plan/\allowbreak cycle/\allowbreak level; current không mất.\\
TC31 /\allowbreak  AT09 & Read không complete & Nhắc Current/\allowbreak Unread có plan chưa làm. & Đánh dấu đọc hai lần; đọc lại plan. & Read idempotent; badge giảm; due/\allowbreak completion không đổi.\\
TC32 /\allowbreak  AT09 & Pause/\allowbreak archive/\allowbreak tắt nhắc & Plan active có nhắc. & Lần lượt pause/\allowbreak archive/\allowbreak disable,  quét; rồi bật/\allowbreak restore/\allowbreak resume. & Không tạo nhắc chủ động khi tắt; đánh giá hiện tại khi bật; không reset kỳ.\\
TC33 /\allowbreak  AT09 & Browser permission denied & Inbox server có một nhắc. & Từ chối quyền browser; mở inbox. & Vẫn đọc được nhắc; không có prompt chặn tác vụ chính.\\
TC34 /\allowbreak  AT10 & Tổng 450.000 và ngày hai đầu & Fuel 20/\allowbreak 10=200.000; items 25/\allowbreak 10=100.000+\allowbreak 150.000. & Report from 20/\allowbreak 10 to 25/\allowbreak 10; so bảng/\allowbreak series/\allowbreak dashboard. & 450.000; hai ngày biên được tính; không cộng header service lần hai.\\
TC35 /\allowbreak  AT10 & Archive filter và data rỗng & Một xe archived có tiền,  một khoảng rỗng. & Đổi includeArchived; query khoảng rỗng và from$>$to. & Mặc định loại archive,  bật thì gồm; rỗng 0; đảo ngày là lỗi.\\
TC36 /\allowbreak  AT11 & Hai tab sửa cùng version & Hai tab tải cùng vehicle hoặc reading v=1. & Tab 1 update thành công; tab 2 update/\allowbreak delete với v=1. & Tab 2 nhận 409; không ghi đè; thông báo tải lại.\\
TC37 /\allowbreak  AT11 & Hai create phá chuỗi liên bản ghi & Fixture riêng: baseline 10.000, ngày 20/\allowbreak 10=10.200 và 25/\allowbreak 10=10.400; hai DB connection đọc cùng history. & Đồng thời tạo 22/\allowbreak 10=10.350 và 23/\allowbreak 10=10.300; mỗi reading hợp lệ nếu xét riêng. & Không thể commit cả hai; guard cấp xe bảo vệ invariant.\\
TC38 /\allowbreak  AT11 & Archive cạnh tranh create nguồn & Hai request cùng xe active. & Đồng thời archive và tạo fuel/\allowbreak service/\allowbreak ODO. & Kết quả tuần tự hợp lệ; không có ghi sau archive đã có hiệu lực.\\
TC39 /\allowbreak  AT12 & Responsive và keyboard & Browser 360/\allowbreak 768/\allowbreak 1440px. & Mở menu,  form,  bảng,  dialog bằng keyboard; kiểm focus/\allowbreak labels. & Không tràn toàn trang; đọc được lỗi; focus không mất; màu có chữ.\\
TC40 /\allowbreak  AT12 & Loading/\allowbreak error/\allowbreak empty/\allowbreak conflict & Network throttling/\allowbreak failure; cache account A. & Mở trang,  gây lỗi,  retry,  logout rồi login B. & Phân biệt states; giữ form/\allowbreak filter; B không thấy dữ liệu A.\\
TC41 /\allowbreak  AT13 & Hiệu năng mục tiêu & 10.000 nguồn; 10 user; warmup và máy đã ghi. & Đo danh sách/\allowbreak report; lưu raw samples và query count. & p95 $\le$1 giây theo NFR05; công bố lỗi/\allowbreak hardware,  không chỉ average.\\
TC42 /\allowbreak  AT13 & Checkout/\allowbreak restart/\allowbreak restore & Checkout sạch; DB demo mới và backup giả. & Cấu hình secrets,  migration; restart; restore DB riêng. & Chạy lại được; count/\allowbreak tổng fixture và ownership giữ đúng.\\
TC43 /\allowbreak  AT14 & Lịch sử không nhân đôi & Fuel/\allowbreak service có reading tự sinh và manual riêng. & Lọc kết hợp,  phân trang với nhiều dòng cùng ngày. & Một activity nguồn chỉ một dòng; manual riêng; sort/\allowbreak paging ổn định.\\
TC44 /\allowbreak  AT15 & Settings owner và validation & A/\allowbreak B có settings; staleOdoDays đích. & A đổi ngưỡng/\allowbreak tắt-bật; gửi 0/\allowbreak âm/\allowbreak lẻ; đọc B. & Không ảnh hưởng B; reject sai; không xóa history; đánh giá hiện tại.\\
TC45 /\allowbreak  AT16 & Cookie/\allowbreak CSRF/\allowbreak session expiry & HTTPS test; cookie/\allowbreak sid hợp lệ. & Bỏ/\allowbreak sai CSRF trên mọi mutation; revoke/\allowbreak expire; kiểm cookie. & CSRF bị từ chối,  cookie HttpOnly/\allowbreak Secure phù hợp,  session cũ 401.\\
TC46 /\allowbreak  AT16 & Đổi password thu hồi phiên & Hai phiên A; password mới hợp policy. & Thử password cũ sai; đổi đúng; dùng lại cả hai phiên. & Sai không đổi; đúng thì cả hai revoked,  yêu cầu login lại theo chính sách tài khoản.\\
TC47 /\allowbreak  AT16 & Admin catalog & User thường và Admin được cấp hợp lệ. & User POST/\allowbreak PUT catalog; Admin làm tương tự; thử stale version. & User 403; Admin được phép khi valid; stale/\allowbreak trùng xử lý conflict.\\
TC48 /\allowbreak  AT17 & DB lỗi và health/\allowbreak docs & Host Development/\allowbreak Production,  DB unavailable. & GET live/\allowbreak ready/\allowbreak docs; gây request DB failure. & Live phản ánh process; ready 503; docs chỉ Development; lỗi không lộ SQL/\allowbreak secret.\\
TC49 /\allowbreak  AT17 & Worker hồi phục và log & Fault injection một vòng worker; dữ liệu giả. & Gây exception,  cho vòng kế tiếp chạy; đọc log/\allowbreak ProblemDetails. & Vòng sau xử lý; traceId/\allowbreak code theo contract đích; không password/\allowbreak token/\allowbreak cookie/\allowbreak connection string.\\
TC50 /\allowbreak AT16 & Thuộc tính cookie phiên & Đăng nhập qua HTTPS; có thể kiểm cookie trong trình duyệt. & 1. Đăng nhập. 2. Kiểm thuộc tính cookie. 3. Thử đọc cookie bằng JavaScript. & Cookie có HttpOnly, Secure và SameSite theo cấu hình; JavaScript không đọc được cookie phiên.\\
TC51 /\allowbreak AT16 & CSRF trên thao tác ghi & A đã đăng nhập; có CSRF token hợp lệ. & 1. Gửi yêu cầu thiếu token. 2. Gửi token sai. 3. Gửi token hợp lệ trên các phân hệ ghi dữ liệu. & Yêu cầu thiếu hoặc sai token bị từ chối, DB không đổi; token hợp lệ cho phép tiếp tục kiểm tra nghiệp vụ.\\
TC52 /\allowbreak AT16 & Khôi phục phiên khi tải lại & A đã đăng nhập; cookie và phiên DB còn hiệu lực. & 1. Tải lại trang. 2. Kiểm /auth/me. 3. Làm hết hạn phiên và tải lại. & Phiên hợp lệ trả hồ sơ A và giữ đăng nhập; phiên hết hạn trả 401 và yêu cầu đăng nhập lại.\\
TC53 /\allowbreak AT16 & Thu hồi phiên trong DB & A có hai phiên độc lập; các mã phiên đã xác định. & 1. Đăng xuất phiên thứ nhất. 2. Dùng lại cookie cũ. 3. Kiểm phiên thứ hai. 4. Đổi mật khẩu và kiểm lại. & Phiên đăng xuất bị thu hồi và trả 401; phiên còn lại giữ hiệu lực đến khi hết hạn hoặc bị thu hồi; đổi mật khẩu thu hồi các phiên.\\
\bottomrule\end{longtable}\end{smalltable}


\section{Kiểm thử phi chức năng}
\subsection{Hiệu năng}
Bộ dữ liệu gồm 10.000 hoạt động và tải đồng thời 10 người dùng. Sau giai đoạn khởi động, nhóm đo thời gian phản hồi của danh sách và báo cáo. Phân vị 95 được đối chiếu với mục tiêu không quá 1 giây; số yêu cầu lỗi được thống kê riêng.

\subsection{Độ trễ nhắc bảo dưỡng}
Nhóm ghi thời điểm kế hoạch thỏa điều kiện nhắc, thời điểm thông báo được lưu và thời điểm xuất hiện trên trang đang mở. Mục tiêu là tạo thông báo trong 120 giây và hiển thị trên trang trong 60 giây tiếp theo khi hệ thống hoạt động bình thường. Các vòng quét lặp và lần khởi động lại phải giữ số lượng thông báo đúng theo khóa kế hoạch, kỳ và mức nhắc.

\subsection{Bảo mật và quản lý phiên}
Theo DEC07, các bài kiểm tra xác thực tập trung vào cookie HttpOnly, Secure, trạng thái phiên trong DB và CSRF. Nhóm kiểm tra yêu cầu thiếu hoặc sai CSRF token, phiên đã hết hạn, phiên đã thu hồi và thao tác đổi mật khẩu. Việc tải lại trang phải giữ đăng nhập khi phiên còn hợp lệ. Khi chuyển từ tài khoản A sang B, giao diện xóa dữ liệu tạm của tài khoản trước.

Quyền truy cập được kiểm tra bằng yêu cầu API trực tiếp. Tài khoản khác không thể đọc, sửa, xóa hoặc liên kết tài nguyên ngoài phạm vi của mình. Nhật ký ứng dụng và phản hồi lỗi không chứa mật khẩu, cookie hoặc thông tin kết nối cơ sở dữ liệu.

\subsection{Khả năng sử dụng}
Nhóm kiểm tra các màn hình tại chiều rộng 360, 768 và 1440 px; thao tác bằng bàn phím; thứ tự focus; nhãn trường và thông báo lỗi. Các tình huống đang tải, chưa có dữ liệu, mất kết nối và xung đột phiên bản cần được phân biệt rõ.

\subsection{Sao lưu và phục hồi}
Trên môi trường kiểm thử, nhóm phục hồi bản sao lưu vào một cơ sở dữ liệu riêng. Kết quả được đối chiếu theo số phương tiện, số hoạt động, tổng tiền và quyền truy cập của hai tài khoản. Thời gian thực hiện được ghi nhận để đánh giá quy trình phục hồi.

\section{Tiêu chí đánh giá}
Các chức năng được chấp nhận khi luồng chính và những ngoại lệ liên quan cho kết quả đúng, không xảy ra truy cập chéo tài khoản, ghi trùng chi phí hoặc lưu dữ liệu một phần. Lỗi mất dữ liệu và sai quyền truy cập được ưu tiên xử lý trước; tiếp theo là lỗi tính hạn, số km, báo cáo và các vấn đề giao diện.

Mỗi lỗi được ghi kèm dữ liệu, bước tái hiện, kết quả mong đợi và kết quả thực tế. Sau khi sửa, nhóm chạy lại ca kiểm thử liên quan và kiểm tra những chức năng chịu ảnh hưởng trước khi tổng hợp kết quả.
